\section{Go Toolchain}
The Go toolchain is the collection of programs and commands used to
develop, build, test, analyse, and package programs written in the Go
programming language.  It is distributed with Go itself and is designed
to provide a standard and largely automated development workflow.\\

The main command is \texttt{go}.  Rather than requiring developers to
invoke a compiler, linker, archiver, test runner, and other tools
individually, the \texttt{go} command coordinates these components.\\

A typical workflow can be represented as

\textbf{
\[
\text{Source code}
\longrightarrow
\text{Packages}
\longrightarrow
\text{Compilation}
\longrightarrow
\text{Linking}
\longrightarrow
\text{Executable}
\]
}

with testing, formatting, dependency management, and analysis occurring
alongside these stages.

\subsection{Installing and Configuring Go}

The first step is to installing a Go distribution.  This provides the
compiler, standard library, \texttt{go} command, and supporting
development tools.\\

After installation, the command is executed.

\begin{lstlisting}
go version
\end{lstlisting}

can be used to verify that Go is available.\\

For example:

\begin{lstlisting}
$ go version
go version go1.25.0 linux/amd64
\end{lstlisting}

The exact version will depend on the installed release.

\subsection*{Why is this needed?}

The compiler and associated tools are required to transform Go source
code into executable programs.  The Go installation also provides the
standard library, which contains commonly used packages such as
\texttt{fmt}, \texttt{os}, \texttt{net/http}, and \texttt{encoding/json}.

\subsection{Creating a Go Module}

Modern Go projects normally use a Go module to define the project's
module path and manage dependencies.\\

A module can be created with this feature.

\begin{lstlisting}
go mod init example.com/hello
\end{lstlisting}

This creates a file named \texttt{go.mod}.  A simplified example is as follows:

\begin{lstlisting}
module example.com/hello

go 1.25
\end{lstlisting}

The module path identifies the project and provides the namespace used
when importing packages belonging to the module.

\subsection*{Why is this needed?}

A Go program may consist of many packages and may depend on packages
provided by other modules.  The module system gives enough information for the Go toolchain  to determine which dependencies are required and
which versions should be used.\\

For example, a project might have the following structure:

\begin{lstlisting}
hello/
|-- go.mod
|-- main.go
|-- internal/
|   `-- greeting/
|       `-- greeting.go
\end{lstlisting}

The \texttt{go.mod} file establishes the root of the module and allows
the toolchain to resolve imports consistently.

\subsection{Writing Go Source Code}

Go programs are normally organised into packages.  The simplest
executable program belongs to the \texttt{main} package and contains a
\texttt{main} function.\\

For example:

\begin{lstlisting}
package main

import "fmt"

func main() {
    fmt.Println("Hello, world!")
}
\end{lstlisting}

The source file can be saved as \texttt{main.go}.

\subsection*{Why is this needed?}

The source code expresses the behaviour of the program.  The Go
toolchain operates on packages rather than treating every source file
as an entirely independent compilation unit.\\

Files belonging to the same package are compiled together.  This allows
the compiler to perform type checking and determine relationships
between declarations within the package.

\subsection{Formatting the Source Code}

Go provides a standard formatter through

\begin{lstlisting}
gofmt
\end{lstlisting}

or, more commonly,

\begin{lstlisting}
go fmt ./...
\end{lstlisting}

For example, code such as

\begin{lstlisting}
func main(){fmt.Println("Hello")}
\end{lstlisting}

will be formatted in the conventional Go style:

\begin{lstlisting}
func main() {
    fmt.Println("Hello")
}
\end{lstlisting}

\subsection*{Why is this needed?}

Formatting is not required for the compiler to understand valid Go
programmes, but it is important for consistency.  Go deliberately has a
strong standard formatting convention.\\

This means that different developers and projects generally use the
same layout, reducing unnecessary arguments about indentation,
braces, spacing, and similar stylistic details.

\subsection{Resolving Dependencies}

If a project imports a package from another module, the Go toolchain
must locate the required module version.

For example:

\begin{lstlisting}
import "example.com/some/library"
\end{lstlisting}

The required dependency is recorded in \texttt{go.mod}.  The command

\begin{lstlisting}
go mod tidy
\end{lstlisting}

can be used to ensure that the module files accurately describe the
dependencies used by the project.\\

The toolchain may also create or update \texttt{go.sum}, which contains
checksums used to verify the downloaded module content.

\subsection*{Why is this needed?}

A compiler cannot compile an import that it cannot locate.  Dependency
resolution therefore establishes the set of packages that are required
before compilation can take place.\\

The module system also makes builds more reproducible by recording
dependency versions rather than relying solely on whatever happens to
be installed on a developer's machine.

\subsection{Loading and Analysing Packages}

When a command such as

\begin{lstlisting}
go build
\end{lstlisting}

When executed, the Go toolchain first determines which packages are
involved and how they depend on one another.\\

For example:

\begin{lstlisting}
main
 |
 +-- internal/greeting
 |
 `-- fmt
\end{lstlisting}

The toolchain constructs the relevant package dependency graph.

\subsection*{Why is this needed?}

The compiler must know which source files belong to each package and
which imported packages are required.  The dependency graph also
allows independent packages to be compiled separately and, where
appropriate, in parallel.

\subsection{Type Checking}

The Go compiler analyses the source code and checks that it obeys the
rules of the Go language.\\

For example:

\begin{lstlisting}
var number int
number = "hello"
\end{lstlisting}

is invalid because a value of type \texttt{string} is being assigned to
a variable of type \texttt{int}.\\

The compiler reports an error rather than producing a normal executable.

\subsection*{Why is this needed?}

Type checking detects many programming errors before the programme is
run.  Go is statically typed, so the compiler can determine many
properties of a program during compilation.

This is one of the reasons for the simplicity of a simple command such as

\begin{lstlisting}
go build
\end{lstlisting}

can catch substantial classes of programming errors.

\subsection{Compiling Packages}

After source analysis and type checking, the Go compiler translates the
source code into compiled object information.\\

Conceptually:

\[
\text{Go source}
\longrightarrow
\text{compiled package}
\]

The compiler is normally invoked indirectly through the \texttt{go}
command rather than being called manually.

For example:

\begin{lstlisting}
go build
\end{lstlisting}

causes the necessary packages to be compiled.

\subsection*{Why is this needed?}

The source code is written for humans and expresses the programme behaviour in the
Go language.  The processor cannot directly execute Go source code.
Compilation translates it into a form suitable for subsequent linking
and execution.\\

The compiler also performs important optimisations and generates
machine code appropriate for the selected target architecture.

\subsection{Linking}

An executable program normally consists of more than the code written
in \texttt{main.go}.  It may depend on many packages, including
packages from the standard library.\\

The linker combines the required compiled components into an executable
program.\\

Conceptually:

\[
\text{compiled packages}
\longrightarrow
\text{linker}
\longrightarrow
\text{executable}
\]

For example:

\begin{lstlisting}
go build -o hello
\end{lstlisting}

can produce an executable named \texttt{hello}.

\subsection*{Why is this needed?}

Compilation produces code for individual packages, whereas an
executable needs to contain the program as a complete linked unit.\\

The linker resolves references between the compiled components and
constructs the final executable image.

\subsection{Building the Program}

The command

\begin{lstlisting}
go build
\end{lstlisting}

coordinates the previous stages.

A simplified model is as follows:

\begin{enumerate}[label=\arabic*.]
    \item Load the module and packages.
    \item Resolve required dependencies.
    \item Compile packages that need compiling.
    \item Link the required components.
    \item Produce an executable when appropriate.
\end{enumerate}

For example:

\begin{lstlisting}
go build -o hello
./hello
\end{lstlisting}

might produce:

\begin{lstlisting}
Hello, world!
\end{lstlisting}

\subsection*{Why is this needed?}

\texttt{go build} provides a convenient abstraction over the compiler
and linker.  Developers normally do not need to know the exact compiler
or linker command required for every package in the program.

\subsection{Running Without Explicitly Building}

Go also provides:

\begin{lstlisting}
go run .
\end{lstlisting}

This builds the necessary program and runs it.

For a small program, one can therefore write:

\begin{lstlisting}
go run main.go
\end{lstlisting}

\subsection*{Why is this needed?}

\texttt{go run} is particularly convenient during development when the
developer wants to execute a program without keeping a separately
managed executable.\\

It should not be interpreted as meaning that the Go source code is directly
executed by the processor.  The necessary compilation still occurs.

\subsection{Testing}

Go has testing support built into its toolchain.  Tests are commonly
placed in files ending in \texttt{\_test.go}.\\

For example:

\begin{lstlisting}
func Add(a, b int) int {
    return a + b
}
\end{lstlisting}

could be tested with:

\begin{lstlisting}
func TestAdd(t *testing.T) {
    if Add(2, 3) != 5 {
        t.Fatal("expected 5")
    }
}
\end{lstlisting}

Tests can be executed using:

\begin{lstlisting}
go test ./...
\end{lstlisting}

Compilation tells us that the program is syntactically and
type-theoretically valid, but it does not prove that the program
produces the desired results.\\

Tests provide executable checks of the program's behaviour. The distinction is important:

\[
\text{Compilation} \neq \text{Correctness}
\]

A program can compile successfully and still contain serious logical
errors.

\subsection{Static Analysis}

Go also provides analysis tools that inspect source code without
executing the program.

A commonly used command is:

\begin{lstlisting}
go vet ./...
\end{lstlisting}

It can identify certain suspicious or incorrect programming constructs.

\subsection*{Why is this needed?}

Some problems are difficult or impossible to identify simply by
checking whether the programme compiles.  Static analysis can identify
patterns that are legal Go but are likely to indicate a programming
mistake.

Thus, a typical quality-orientated workflow is as follows.

\[
\text{format}
\rightarrow
\text{build}
\rightarrow
\text{test}
\rightarrow
\text{analyse}
\]

\subsection{Cross Compilation}

One notable feature of the Go toolchain is straightforward
cross-compilation.\\

For example, a developer on Linux can build for another operating
system and architecture by setting environment variables such as
\texttt{GOOS} and \texttt{GOARCH}:

\begin{lstlisting}
GOOS=windows GOARCH=amd64 go build -o hello.exe
\end{lstlisting}

The resulting executable targets Windows on the AMD64 architecture.

\subsection*{Why is this needed?}

Software is often developed on one platform but deployed on another.
Cross-compilation allows a developer or build server to produce
binaries for deployment targets without necessarily running the build
on the target operating system.\\

Conceptually:

\[
\text{Development machine}
\xrightarrow{\text{Go toolchain}}
\text{Target executable}
\]

\subsection{Caching}

The Go toolchain maintains build caches.  When a package has already
been compiled and its relevant inputs have not changed, the toolchain
can reuse the cached result.\\

Large projects can contain many packages.  Recompiling every package
after every small source-code change would be unnecessarily expensive.\\

Caching therefore reduces build times while preserving correctness:
if an input that affects a package changes, the relevant package is
rebuilt.\\

\subsection{The Complete Development Cycle}

A practical Go development workflow can therefore be summarised as:

\begin{enumerate}
    \item \textbf{Initialise a module} using \texttt{go mod init}.
    \item \textbf{Write source code} organised into packages.
    \item \textbf{Format the code} using \texttt{gofmt} or
          \texttt{go fmt}.
    \item \textbf{Resolve dependencies} using the Go module system.
    \item \textbf{Compile and type-check} the packages.
    \item \textbf{Link} the compiled packages into an executable.
    \item \textbf{Run} the program using \texttt{go run} or the
          resulting executable.
    \item \textbf{Test} the program using \texttt{go test}.
    \item \textbf{Analyse} it using tools such as \texttt{go vet}.
    \item \textbf{Build for deployment}, potentially using
          cross-compilation.
\end{enumerate}

The relationship between these activities can be visualised as:

\[
\boxed{
\begin{array}{c}
\text{Go source files}\\
\downarrow\\
\text{Formatting}\\
\downarrow\\
\text{Package loading}\\
\downarrow\\
\text{Dependency resolution}\\
\downarrow\\
\text{Parsing and type checking}\\
\downarrow\\
\text{Compilation}\\
\downarrow\\
\text{Linking}\\
\downarrow\\
\text{Executable}
\end{array}}
\]

Testing and static analysis operate alongside this pipeline:

\[
\begin{array}{ccccc}
& \text{Source code} & & &\\
& \downarrow & & &\\
\text{Formatting} & \rightarrow & \text{Build} & \rightarrow &
\text{Executable}\\
& & \downarrow & &\\
& & \text{Tests} & &\\
& & \downarrow & &\\
& & \text{Static analysis} & &
\end{array}
\]

\subsection{Conclusion}

The Go toolchain is deliberately integrated.  The developer generally
interacts with the single \texttt{go} command rather than manually
coordinating separate compilation, linking, dependency-management,
testing, and packaging commands.\\

Its central purpose is to turn a collection of Go source files and
their dependencies into a reliable, executable programme while making
common development activities predictable and repeatable.\\

In simplified form:

\[
\boxed{
\text{Source}
\rightarrow
\text{Packages}
\rightarrow
\text{Type checking}
\rightarrow
\text{Compilation}
\rightarrow
\text{Linking}
\rightarrow
\text{Executable}
}
\]

with

\[
\boxed{
\text{Formatting}
\quad
\text{Testing}
\quad
\text{Static analysis}
\quad
\text{Dependency management}
}
\]

supporting the main build pipeline.